\section{ROAR：3D 重光照}

\subsection{动机：细粒度控制}

除了生成3D内容本身，对输出进行细粒度控制也至关重要——\textbf{光照}是其中最重要的控制维度之一。ROAR 方法解决的问题是：给定一组带位姿的图像和一个目标光照条件，生成该光照下的3D模型。

\subsection{方法设计}

ROAR 训练了一个\textbf{多视角重光照模型}，其架构类似于 CAT3D 和 B3D 中的多视角模型，但增加了光照条件作为输入。

具体而言：
\begin{itemize}
    \item 输入：一组图像、对应位姿 + 目标光照条件
    \item 输出：每个视角在目标光照下的重光照图像
    \item 目标：对于 $N$ 个视角和 $M$ 个光照条件，生成 $N \times M$ 张重光照图像
    \item 这些图像随后用于训练一个\textbf{光照条件 NeRF}
\end{itemize}

核心思路类似虚拟光舞台（light stage）——通过模拟不同光照条件下的物体外观，训练 NeRF 对任意目标光照进行泛化。

\subsection{光照表示：环境贴图 + 镜面反射嵌入}

光照的编码是一个非平凡的问题。ROAR 采用了双重光照表示：

\begin{importantbox}{ROAR 的双重光照编码}
\textbf{1. 全局光照嵌入}：将环境贴图（environment map）通过 Transformer 编码器转换为光照嵌入向量。

\textbf{2. 镜面反射嵌入}：由于全局编码器不够强大，无法完全捕获镜面高光细节，ROAR 额外引入了一个镜面反射嵌入（specular embedding），从镜面反射方向对环境贴图进行局部采样。为了处理不同粗糙度，采用预模糊的环境贴图，从而只需采样单个位置即可获得镜面反射信息。
\end{importantbox}

\subsection{实验效果}

实验表明，ROAR 能够忠实地恢复不同目标光照下的镜面高光效果，而竞争方法要么遗漏镜面高光，要么过度补偿。

\subsection{本章小结}

ROAR 将「2D生成 + 3D蒸馏」范式扩展到了光照控制维度，通过训练多视角重光照模型生成丰富的光照-视角数据，再蒸馏为光照条件 NeRF，实现了3D物体的任意重光照。

